\documentclass[10pt, a4paper]{article}
\usepackage[T2A]{fontenc}
\usepackage[russian]{babel}
\usepackage{geometry}
\usepackage{titlesec}
\usepackage{tikz}
\usepackage{amsmath}
\usepackage{amssymb}
\usepackage{array}

\setlength{\parindent}{0pt} % Убираем абзацные отступы

\titleformat{\section}[block]{\centering\Large\bfseries}{}{0pt}{}
\titleformat{\subsection}[block]{\centering\large\bfseries}{}{0pt}{}
\setcounter{secnumdepth}{-1}

\geometry{top=2cm, bottom=2cm, left=2cm, right=2cm}

\title{
    \textbf{Федеральное государственное автономное образовательное учреждение высшего образования}\\
    \textbf{«Национальный исследовательский университет}\\
    \textbf{"Высшая школа экономики"»}

    \vspace{1cm}

    \textbf{Московский институт электроники и математики им. А.Н. Тихонова НИУ ВШЭ}

    \textbf{Департамент компьютерной инженерии}

    \vspace{1.5cm}

    \textbf{Курс: Алгоритмизация и программирование}

    \vspace{1cm}

    \textbf{ОТЧЕТ}

    \textbf{по лабораторной работе №1}
}

\author{
    \fontsize{10}{10}
    \begin{tabular}{cc}
        \begin{tabular}{|l|c|c|}
            \hline
            \textbf{Раздел} & \textbf{Mах оценка} & \textbf{Итог. оценка} \\
            \hline
            Постановка & $0,5$ & \\
            \hline
            Метод & $1$ & \\
            \hline
            Спецификация & $0,5$ & \\
            \hline
            Алгоритм & $1,5$ & \\
            \hline
            Работа программы & $1$ & \\
            \hline
            Листинг & $0,5$ & \\
            \hline
            Тесты & $1$ & \\
            \hline
            Вопросы & $2$ & \\
            \hline
            Доп. задание & 2 & \\
            \hline
        \end{tabular}
        &
        \begin{tabular}{l}
            Студент: Солодилов Михаил Анатольевич \\
            Группа: БИВ253 \\
            Вариант: 232\\
            Руководитель:\\
            Оценка:\\
            Дата сдачи \\
        \end{tabular}
    \end{tabular}
}

\date{
    \vfill
    \textbf{МОСКВА 2025}
}

\addto\captionsrussian{\renewcommand\contentsname{Оглавление}}

\begin{document}

\maketitle
\thispagestyle{empty}
\newpage
\pagenumbering{arabic}
\tableofcontents
\newpage

\section{Задание}

\begin{enumerate}
    \item (2) Вычислить массив $R[1:n]$ в соответствии с формулой:
        $r[i] = 1,25 \sin(3ax -ih)$
    \item (8) Из вычисленного массива $R$ удалить все элементны,
        расположенные между последним положительным и первым
        максимальным по модулю элементами.
    \item (3) В полученном массиве $R[1:k]$, где $k$ --- число
        элементов, оставшихся после удаления, подсчитать среднее
        арифметическое элементов, расположенных до последнего
        минимального элемента включительно.
\end{enumerate}

\section{Постановка задачи}

\underline{Дано:}
\begin{enumerate}
    \item n --- цел.; a, x, h --- вещ.
    \item Нет входных данных
    \item Нет входных данных
\end{enumerate}

\underline{Результат:}
\begin{enumerate}
    \item $R[1:n]$ --- вещ.
    \item $R[1:k]$ --- вещ., где $k$ --- длина массива после удаления
        или сообщение <<Full removal>> или сообщение <<No deletions>>
    \item $\text{average}$ --- вещ. или сообщение <<No average>>
\end{enumerate}

\underline{При:}
\begin{enumerate}
    \item $n \in \mathbb{N}, \ n \leq LMAX$, где $LMAX$ ---
        максимальная длинна массива.
\end{enumerate}

\underline{Связь:}
\begin{enumerate}
    \item $r[i] = 1,25 \sin(3ax - ih)$
    \item $\exists n_1 : \forall i \in [1, n] \implies |r[i]|
        <= |r[n_1]|, \  \nexists t \in [1, n_1 - 1]: |r[t]| >
        |r[n_1]|$ \\
        $\exists n_2 : r[n_2] > 0 \  \nexists t \in [n_2 + 1, n]:
        r[t] > 0$ \\
        При $n_1 > n_2, \ c = n_1, \ n_1 = n_2, \ n_2 = c$ \\
        $\forall i \in [1, n_1] \cup [n_2, n] \implies \exists t \in
        [1, k]: r'[t] = r[k], \ \forall j \in [1, k] \ \exists q \in
        [1, n_1] \cup [n_2, n]: r'[j] = r[q]$, где $r'$ --- массив
        после удаления.
    \item $\exists km \in [1, k] : \forall i \in [1, k] \implies r[i]
        >= r[km], \ \forall q \in [km + 1, k] \implies r[q] > r[km]$
        \\
        $average = \sum_{i = 1}^{i = km}{r[i]} \cdot \frac{1}{km}$
\end{enumerate}

\section{Метод решения задачи}
\begin{enumerate}
    \item
        $
            \begin{array}{l}
                \begin{cases}
                    \text{для } i = \overline{1, n} \\
                    r[i] = 1,25 \sin(3ax - ih)
                \end{cases}
            \end{array}
        $
    \item
        $
            \begin{array}{l}
                n_1 = 0 \\
                \begin{cases}
                    для i = \overline{2, n} \\
                    n_1 = i \text{, если } |r[i]| > |r[n_1]| \\
                \end{cases} \\ \\

                n_2 = n \\
                \begin{cases}
                    \text{пока } n_2 > 0 \text{ и } r[n_2] \leq 0 \\
                    n_2 = n_2 - 1 \\
                \end{cases} \\ \\

                m = n_2 - n_1  - 1\\
                \begin{cases}
                    \text{для } i = \overline{n_1 + 1, n_2 - 1} \\
                    r[i] = r[i + m] \\
                \end{cases} \\
                k = n - m \\
            \end{array}
        $
    \item
        $
            \begin{array}{l}
                km = 0
                \begin{cases}
                    \text{для } i = \overline{2, k} \\
                    km = i \text{, если } r[i] \leq r[km]
                \end{cases} \\ \\

                sum = 0 \\
                \begin{cases}
                    \text{для } i = \overline{1, km} \\
                    sum = sum + r[i] \\
                \end{cases} \\
                average = \frac{sum}{km}
            \end{array}
        $
\end{enumerate}

\section{Внешняя спецификация}



\section{Описание алгоритма на псевдокоде}
\section{Листинг программы}
\section{Распечатка тестов к программе и результатов}


\end{document}
